\section{Background}
\label{sec:rs-background}

The Recursive Spawning Protocol (RSP) operates at the intersection of several established research and engineering traditions: agent lifecycle management in multi-agent AI systems, accountability and provenance in distributed autonomous systems, fixed versus mutable delegation structures, hierarchical task decomposition in classical and modern AI, and the BX3 Framework's three-layer model as the architectural substrate for immutable parent chains. Each tradition contributes an essential conceptual component; RSP's contribution is integrating them into a unified structural solution to the accountability gap that emerges when agents spawn children autonomously.

% -------------------------------------------------------
\subsection{Agent Lifecycle Management in Multi-Agent Systems}
% -------------------------------------------------------

Multi-agent systems (MAS) research has long recognized that an agent's lifecycle --- its creation, activation, operation, suspension, and termination --- is a first-class design concern, not merely an implementation detail. The foundational taxonomy introduced by Stone and Veloso \cite{stone2000} classifies multi-agent systems along axes of agent heterogeneity (homogeneous versus heterogeneous populations), agent knowledge (perfect versus imperfect information), environment accessibility (fully versus partially accessible), and the nature of agent interactions (cooperative versus competitive). Within this taxonomy, recursive spawning is most relevant to cooperative, homogeneous systems operating in partially accessible environments --- precisely the regime in which modern agentic runtimes such as AutoGPT \cite{autoGPT2023}, LangChain \cite{langchain2023}, and Microsoft's AutoGen \cite{autogen2023} operate.

Classical agent lifecycle management was primarily an engineering concern: spawning, scheduling, and terminating agents to maximize throughput or minimize latency, with accountability residing with the system operator who configured the agent population \cite{bansal2019}. As AI agents move into healthcare, infrastructure control, and financial decision-making, this assumption becomes untenable. When a spawned agent's action produces real-world consequences, the question of who authorized the spawn, through what chain of delegation, and with what scope constraints is no longer academic --- it is a legal, regulatory, and safety question with no clean answer in conventional architectures.

Ferber and M\"uller \cite{ferber1996} formalized lifecycle interactions as a state machine coupled to the organizational context in which an agent operates, demonstrating that organizational roles --- not just individual agent states --- must be modeled to ensure system-wide coherence. Padgham and Thiel \cite{padgham2001} further established that incomplete lifecycle management, particularly the failure to distinguish between suspension and termination states, introduces systemic failures where orphaned or zombie agents continue consuming resources and emitting effects after their parent has exited.

Recent work has begun addressing the accountability gap in multi-agent lifecycles directly. The Human Delegation Provenance (HDP) protocol identifies the core vulnerability in multi-hop delegations: in chains of the form human $\rightarrow$ orchestrator $\rightarrow$ sub-agent $\rightarrow$ tool-agent, verifying that terminal actions were genuinely authorized by a human requires reconstructing the full delegation chain --- a task current standards (OAuth 2.0, JWTs, UCAN) do not adequately address for append-only, human-provenance requirements in agentic systems \cite{hdp-draft,hdp-draft-2}. HDP proposes cryptographic tokens capturing each delegation hop as a signed entry in an append-only chain, enabling offline verification using only the issuer's public key and session ID. The Agent Lifecycle Protocol (ALP) extends this with canonical lifecycle events --- Genesis, Fork, Migration, Retirement, Succession, and Decommission --- with explicit state-machine semantics, transition rules, and a genealogical registry capturing both genetic (model, architecture) and epigenetic (config, memory, reputation history) lineage for complete provenance across delegation chains \cite{alp2025}. PROV-AGENT demonstrates convergence with the W3C PROV standard by extending it with the Model Context Protocol to link agent decisions and prompts with end-to-end workflow data, enabling near real-time provenance capture across facility boundaries \cite{prov-agent2025}. Unified Agent Lifecycle Management (UALM) provides a complementary five-layer blueprint for end-to-end governance of agent fleets: identity and persona registry, orchestration and cross-domain mediation, PHI-bounded context and memory, runtime policy enforcement with kill-switch triggers, and lifecycle management linked to credential revocation and audit logging \cite{ualm2025}.

Together, these protocols establish that the industry has identified the accountability gap in agentic lifecycles and is converging on cryptographic, append-only delegation records as the foundational solution. RSP's contribution relative to these protocols is the structural enforcement layer: making these delegation records immutable by architectural construction rather than by operational convention.

The core problem RSP addresses is \emph{spawn chain opacity}: the inability to determine, at any point during or after execution, whether a given agent's authority traces to a human principal or to an autonomous predecessor. This opacity arises because parent pointers are mutable in conventional architectures: they can be redirected by an administrative action, a policy override, or a runtime reconfiguration. If the parent pointer of an agent can be redirected after provisioning, then the agent's current parent may not be the principal who authorized its spawn, and the chain of accountability is incomplete or misleading. RSP eliminates this ambiguity by making the parent pointer architecturally immutable from the moment of spawn.

% -------------------------------------------------------
\subsection{Accountability and Provenance in Distributed Systems}
% -------------------------------------------------------

Provenance --- the record of how a data item or decision came to exist --- is a well-studied problem in distributed systems and scientific data management. The W3C PROV data model \cite{moreau2015-prov} formalizes provenance as a set of entities, activities, and agents linked by relations of derivation, attribution, and responsibility. Under PROV-DM, any entity can be attributed to an agent, and agents can themselves be attributed to other agents, enabling the construction of a provenance graph encoding the full chain of responsibility for a given artifact. The core insight PROV contributes to RSP is that the provenance graph is only as reliable as the immutability guarantees of its edges: if the graph is mutable, any attribution derived from it is conditional rather than definitive.

The immutable ledger tradition in distributed systems provides the structural solution. Certificate Transparency \cite{laurie2013certificate} demonstrates that append-only logging with hash-chaining can produce tamper-evident records without centralized trust: any auditor with read access to the log can independently verify that no entry has been inserted, deleted, or modified by checking the hash chain. Bitcoin's blockchain \cite{nakamoto2008bitcoin} extends this with a consensus-based append-only ledger, demonstrating that sufficiently distributed consensus can produce globally agreed records with no trusted central authority. Hyperledger Fabric \cite{hyperledger2015} adapts this model to permissioned contexts, showing that hash-chained ledgers are deployable in governed environments where not all participants are mutually trusting.

Applied to agentic systems, these patterns establish the feasibility of a forensic ledger: an append-only, tamper-evident record of every spawn event, state transition, and delegation action, maintained by the runtime infrastructure rather than by individual agents. Such a ledger enables retrospective accountability: given any agent's identity, a complete chain of delegations can be reconstructed; given any decision or action, the full set of principals who authorized it can be identified. The ledger does not require centralized trust because any observer with read access can independently verify its integrity via the hash chain.

The accountability challenge in distributed agentic systems is distinctive because agents are not passive data artifacts --- they are autonomous actors capable of spawning children, issuing directives, and modifying their own behavior at runtime. This makes the standard provenance model insufficient: the provenance graph of a multi-agent decision may span dozens of agents and hundreds of spawn events, and it must be traversable in real time to support the upstream accountability guarantee (Section~\ref{sec:rs-integration}). A forensic ledger for agentic systems therefore requires not only append-only storage and hash-chaining but also low-latency write performance, indexed retrieval, and protocol-level write enforcement --- properties that the Fact Layer of the BX3 Framework is specifically designed to provide.

% -------------------------------------------------------
\subsection{Fixed Versus Mutable Delegation Chains}
% -------------------------------------------------------

Delegation chains in multi-agent systems take two structurally distinct forms with profoundly different accountability properties. In a \emph{fixed delegation chain}, the parent-child relationship is established at spawn time and is architecturally immutable thereafter: the child's parent pointer cannot be changed, the child's authority cannot be re-delegated except through the defined spawn protocol, and the chain terminates at a predetermined anchor (typically a human principal). In a \emph{mutable delegation chain}, the parent-child relationship is treated as runtime state that can be updated, redirected, or reassigned as conditions change.

Mutable delegation chains offer operational flexibility: a child agent can be re-parented to a different parent if the organizational structure changes, if a parent fails, or if load balancing requires reassignment. This flexibility, however, comes at a significant accountability cost. If the parent-child relationship is mutable at runtime, then any single action by a child agent may have been authorized by one parent at one moment and a different parent at another --- and the effective authorization for a given action may not be the authorization that was recorded at spawn time. The chain of provenance becomes conditional on the state of a mutable mapping rather than fixed by an immutable record. This introduces the \emph{delegatee ambiguity problem}: after an event, it may be impossible to determine which principal effectively authorized the action, because the delegation chain has been rewritten since the action occurred.

The argument for mutability is operational necessity: in dynamic environments, task requirements change, agents fail, and organizational structures evolve. A static parent pointer becomes a liability when the parent is decommissioned. The conventional response is governance by policy: administrative documents specify conditions under which parent pointers may be redirected, and audit logs record redirection events. This approach works in low-stakes environments but fails in high-stakes ones where the cost of an unauthorized redirection may be catastrophic and irreversible. The HDP protocol identifies this precisely as the core vulnerability in conventional delegation: without cryptographic chain-of-custody, there is no verifiable link between a terminal action and its original human authorizer \cite{hdp-draft,hdp-draft-2}.

A fixed delegation chain permits no modification to the parent pointer after provisioning. The chain is established once, at spawn time, and persists for the agent's lifetime. This is a stronger constraint than access-control policy: the Fact Layer enforces immutability of the parent pointer through write protocol, rejecting any request to overwrite $\alpha(n)$ after provisioning regardless of who issued the request or what justification was offered. The immutability is structural, not contractual. The fixed chain does not preclude dynamic task reassignment --- it requires that reassignment be implemented through the spawn protocol rather than retroactive pointer modification. When a task must be reassigned, the agent terminates and a new agent is spawned under the appropriate parent. The distinction matters because the spawn event is a Fact Layer write, creating a new provenance record with full metadata, whereas a pointer modification would retroactively obscure the original authorization.

The tradeoff corresponds to what the distributed systems literature terms \emph{strong consistency} versus \emph{eventual consistency} \cite{vogels2009eventual}: fixed chains trade the ability to reparent dynamically for the ability to reason precisely about authorization state at all times. RSP adopts the fixed chain model because the accountability requirements of autonomous agentic systems --- particularly in safety-critical domains such as irrigation control, infrastructure management, and clinical decision support --- demand strong, unconditionally verifiable accountability chains rather than eventually consistent approximations.

% -------------------------------------------------------
\subsection{Hierarchical Task Decomposition in AI}
% -------------------------------------------------------

Hierarchical task decomposition (HTD) is one of the foundational techniques in classical AI planning. Hierarchical Task Network (HTN) planning, introduced by Ghallab et al. \cite{ghallab2004htn} and formalized in \cite{erpl}, decomposes a high-level task into a hierarchy of subtasks, where each decomposition is a method that specifies how a compound task can be replaced by a partially ordered set of subtasks. The decomposition terminates when all primitive tasks are executable by the system's primitive action set. HTN planning's key properties are that decomposition is pre-planned (the full decomposition tree is constructed before execution begins) and that the hierarchy is explicitly represented in the planning domain as a set of task schemata.

Simon's concept of bounded rationality \cite{simon} provides the theoretical foundation for why hierarchical decomposition is necessary: due to cognitive limits, agents cannot evaluate all possible courses of action at once and must decompose complex problems into manageable subproblems. This insight applies directly to modern LLM-based agents: an agent with a complex goal must decompose it because the full goal space exceeds its representational and computational capacity.

Modern LLM-based agents have revived hierarchical task decomposition under the rubric of dynamic, runtime spawning. Unlike classical HTN planning where decomposition is pre-planned and static, LLM-based agents decompose tasks dynamically in response to encountered subproblems: an LLM agent can observe that a current action has failed, infer that the failure stems from task complexity beyond its provisioned scope, and autonomously spawn a child agent to handle a specific sub-task \cite{press2022,yao2022}. The THREAD system \cite{thread2025} provides perhaps the closest conceptual relative to RSP in this tradition, treating LLM generation as a dynamic multi-threaded process where a parent thread can spawn child threads to handle sub-tasks (think, retrieve data, plan steps) and pause until those children return concise results. This dynamic recursive spawning enables hierarchical task decomposition while managing context length.

What connects classical HTN and modern LLM-based decomposition is the spawn tree: a hierarchical structure in which each node represents a subtask and edges represent the decomposition relationship. In both traditions, the tree is typically mutable and implicit --- decomposition decisions are not recorded as first-class authorization artifacts, and child nodes are not distinguished from parent nodes in terms of their accountability status. A classical HTN planner's subtasks are subroutine calls; they do not have independent accountability chains. An LLM agent's spawned sub-agents are independent processes, but the chain of authorization is implicit in the agent's context window rather than explicit in an immutable warrant structure. Dynamic HTD in open-loop agentic systems has been identified as a vector for \emph{autonomous drift}: the progressive expansion of agent behavior beyond authorized scope, precisely because decomposition is unconstrained \cite{autoGPT2023}.

Recursive Spawning adopts the hierarchical decomposition paradigm from both traditions but adds the accountability layer that is absent from both: each node in the decomposition tree is a first-class accountability entity with an immutable parent pointer, a defined termination condition, and a cryptographically verifiable chain back to a human anchor. Decomposition is not merely a planning technique; it is an authorization event that creates a new accountable entity. The decomposition tree is therefore not just a planning artifact but an organizational chart of delegated authority.

% -------------------------------------------------------
\subsection{The BX3 Framework as Substrate for Immutable Parent Chains}
% -------------------------------------------------------

The BX3 Framework \cite{bx3-framework,bx3framework2026} provides the architectural substrate within which Recursive Spawning operates. The BX3 Framework organizes any autonomous node into three immutable functional layers --- the \textit{Purpose Layer}, the \textit{Bounds Engine}, and the \textit{Fact Layer} --- each defined by the properties it must maintain rather than by the type of actor occupying it. The Purpose Layer carries intent, judgment, and final human accountability; the Bounds Engine carries constrained execution within a defined operating range; the Fact Layer carries deterministic enforcement and forensic auditability.

The upstream accountability guarantee of the BX3 Framework states that when any node encounters a state it cannot resolve within its provisioned bounds, accountability escalates recursively upward through the system hierarchy, bypassing all machine actors, until it reaches a human anchor \cite{bx3-framework,bx3framework2026}. This guarantee is not a policy choice; it is a structural property of the three-layer architecture. Specifically, the Fact Layer creates the immutable parent pointer at spawn time, issues the spawn warrant, and maintains the forensic ledger that records every spawn event. The Fact Layer's determinism is what makes the parent pointer architecturally immutable: once written, the pointer cannot be updated because the Fact Layer does not permit modifications to its authorization records. Any attempt to overwrite $\alpha(n)$ is rejected at the protocol level before execution. This contrasts with ACL-based immutability in conventional systems: administrative compromise in ACL-based systems leads directly to immutability violation, whereas BX3's protocol-level enforcement has no administrative override path.

The depth bound $D_{\max}$ is a Purpose-Layer-defined, Bounds-Engine-enforced, Fact-Layer-hardened constraint. The Purpose Layer defines the maximum permissible chain depth as part of the spawn authorization policy. The Bounds Engine evaluates each proposed spawn against this limit, returning \texttt{SPAWN\_REFUSED} if the spawn would exceed it. The Fact Layer's pre-commit hook enforces this limit structurally: any spawn operation that would produce $|C| \geq D_{\max}$ is rejected before execution, regardless of the Bounds Engine's return. The constraint is thus enforced three times --- by the layer that defines it, the layer that evaluates against it, and the layer that structurally blocks violations.

The forensic ledger is a Fact Layer artifact. Every spawn event, state transition, Purpose Layer directive, and exception condition is recorded in the append-only ledger maintained by the Fact Layer. Each entry is hash-chained to its predecessor, making the ledger tamper-evident: any insertion, deletion, or modification of a ledger entry breaks the hash chain and is detectable by any observer with read access. The ledger provides the evidentiary foundation for RSP's correctness properties (Section~\ref{sec:rs-guarantees}) because it contains a complete, immutable record of every decision point in the chain's lifecycle.

The human anchor $h(n)$ is a Purpose Layer assignment. At provisioning, the Purpose Layer records the identity of the human principal who authorized the root spawn event. This assignment is immutable for the lifetime of the chain: for all nodes $n$ in chain $C$, $h(n) = h(n_0)$ regardless of chain depth. This is the architectural expression of the upstream BX3 accountability guarantee: systems built on the BX3 Framework fail upward into human consciousness, never downward into autonomous chaos.

The structural necessity of the three-layer mapping is established formally in Section~\ref{sec:rs-integration}. Removing any layer causes the corresponding guarantee to fail. Without the Purpose Layer, there is no authoritative human anchor and no source of spawn authorization policy. Without the Bounds Engine, depth limits become unenforceable conventions subject to override by operational urgency. Without the Fact Layer, immutability of the parent pointer and tamper-evidence of the ledger are promises rather than constraints. The three-layer architecture is not a design pattern --- it is the minimal structure that makes RSP's correctness properties unconditionally guaranteed rather than conditionally enforced.
